En réponse à nos questions, notre \robot\ peut indiquer dans quelle direction il est orienté.
N'étant pas doué de la parole, ses réponses sont {\em booléennes} \index{Booléen} c'est-à-dire 
{\idt true} (prononcer ``trou'', {\em vrai} \index{Vrai} en anglais) ou {\idt false} (prononcer ``faulse'', 
{\em faux} \index{Faux} en anglais). On peut imaginer qu'il donne ses réponses à l'aide d'un voyant lumineux
qu'il allume pour répondre  {\idt true} (vrai) et qu'il laisse éteint pour répondre {\idt false}
(faux).

Pour poser les questions sur l'orientation du \robot, on utilise les commandes de test suivantes~:

\begin{table}[htbp]
\begin{center}
\begin{tabular}{@{}|l|p{11.7cm}|@{}} \hline
\multicolumn{2}{|l|}{\em Commandes de test d'orientation du \robot~:} \\ \hline \hline
{\idt tn} & {\idt true} si  le \robot\ est \underline{t}ourné vers le nord (\underline{n}) et  {\idt false} sinon. \\ \hline
{\idt tne} & {\idt true} si  le \robot\ est \underline{t}ourné vers le nord-est (\underline{ne}) et {\idt false} sinon. \\ \hline
{\idt te} & {\idt true} si  le \robot\ est \underline{t}ourné vers l'est (\underline{e}) et {\idt false} sinon. \\ \hline
{\idt tse} & {\idt true} si  le \robot\ est \underline{t}ourné vers le sud-est (\underline{se}) et {\idt false} sinon. \\ \hline
{\idt ts} & {\idt true} si  le \robot\ est \underline{t}ourné vers le sud (\underline{s}) et {\idt false} sinon. \\ \hline
{\idt tsw} & {\idt true} si  le \robot\ est \underline{t}ourné vers le sud-ouest (\underline{sw}) et {\idt false} sinon. \\ \hline
{\idt tw} & {\idt true} si  le \robot\ est \underline{t}ourné vers l' ouest (\underline{w}) et {\idt false} sinon. \\ \hline
{\idt tnw} & {\idt true} si  le \robot\ est \underline{t}ourné vers le nord-ouest (\underline{nw}) et {\idt false} sinon. \\ \hline
\end{tabular}
\end{center}
%\label{tab:tnetc}
\caption{Tests d'orientation du \protect\robot}
\end{table}

En \pascal, l'instruction de test ``if'' permet de n'exécuter une séquence d'instructions que
si une certaine {\em expression booléenne} \index{Expression booléenne}est vraie ({\idt true}).  Pour commencer nous allons
utiliser des expressions booléennes simples, à savoir les commandes de test d'orientation du \robot\
({\idt tn}, {\idt tne}, {\idt te}, {\idt tse}, {\idt ts}, {\idt tsw}, {\idt tw}, {\idt tnw}). Dans sa
forme la plus simple, l'instruction de test ``if'' s'écrit comme suit~:

\index{Test}
\begin{center}
\begin{tabular}{@{}|p{12.6cm}|@{}} \hline
{\em Test ``if'' à une seule alternative~:}\\ \hline \hline
\begin{minipage}[t]{12.6cm}
{\medskip
\mbox{\blc\ \ \ {\mcl if}\hspace{0.3em}\blanc\hspace{0.3em}{\small\rm expression booléenne}%
\footnote{\ A résultat {\idt true} (vrai) ou {\idt false} (faux).} \hspace{0.3em}\blanc\hspace{0.3em}{\mcl then}}\par 
\mbox{\blc\ \ \ \ \mcl begin}\par 
\mbox{\blc\ \ \ \ \ \small\rm  \ldots\ Instructions (séparées par des `$;$') exécutées seulement si}\par 
\mbox{\blc\ \ \ \ \ \small\rm  l'expression booléenne vaut {\idt true} \ldots}\par 
\mbox{\blc\ \ \ \ {\mcl end}\idt$;$}\par 
\medskip}
\end{minipage}\\ \hline 
\end{tabular}
\end{center}

On peut traduire en français par {\idt si} \ldots\ {\idt alors} {\idt début}  \ldots\ {\idt fin}.
Le mot anglais {\idt then} se prononce très approximativement `zen' (avec un cheveu sur la langue).

\underline{Une erreur grave est de mettre un point-virgule 
`{\idt $;$}' après le {\mcl then}}. \hfill Dans ce cas les instructions figurant entre le {\mcl begin} et
le {\mcl end} sont toujours exé\-cu\-tées. 

Dans l'exemple ci-dessous, le \robot\ avance deux fois et tourne à droite d'un quart de tour puis
d'un quart de tour supplémentaire 
s'il est orienté à l'ouest~:

\vspace{\baselineskip}

 {\raggedright 
 \includevspacePostScript[0.9cm]{FriseEnPeigne}{0.95}{0.95}
 \input FriseEnPeigne.tex
}
 
 Dans sa première forme ci-dessus l'instruction de test ne fait rien quand l'expression booléenne
 vaut {\idt false}. Si l'on veut donner le choix de la séquence d'instructions à exécuter dans les
 deux alternatives, on utilisera la forme suivante~:
 
\index{Test}
 \begin{center}
\begin{tabular}{@{}|p{12.6cm}|@{}} \hline
{\em Test ``if'' à deux alternatives~:}\\ \hline \hline
\begin{minipage}[t]{12.6cm}
{\medskip
\mbox{\blc\ \ \ {\mcl if}\hspace{0.3em}\blanc\hspace{0.3em}{\small\rm expression booléenne}%
\footnote{ \ A résultat {\idt true} (vrai) ou {\idt false} (faux).} \hspace{0.3em}\blanc\hspace{0.3em}{\mcl then}}\par 
\mbox{\blc\ \ \ \ \mcl begin}\par 
\mbox{\blc\ \ \ \ \ \small\rm  \ldots\ Instructions (séparées par des `$;$') exécutées si l'expression}\par 
\mbox{\blc\ \ \ \ \ \small\rm   booléenne vaut {\idt true} \ldots}\par 
\mbox{\blc\ \ \ \ {\mcl end}}\par 
\mbox{\blc\ \ \ {\mcl else}}\par 
\mbox{\blc\ \ \ \ \mcl begin}\par 
\mbox{\blc\ \ \ \ \ \small\rm  \ldots\ Instructions  (séparées par des `$;$') exécutées si l'expression  }\par 
\mbox{\blc\ \ \ \ \ \small\rm  booléenne vaut {\idt false} \ldots}\par 
\mbox{\blc\ \ \ \ {\mcl end}\idt$;$}\par 
\medskip}
\end{minipage}\\ \hline 
\end{tabular}
\end{center}

On peut traduire en français par {\idt si} \ldots\ {\idt alors} {\idt début}  \ldots\ {\idt fin}
 \ldots\ {\idt sinon} {\idt début}  \ldots\ {\idt fin}.

\underline{Une erreur grave est de mettre un point-virgule 
`{\idt $;$}' après le {\mcl else}}. \hfill Dans ce cas les instructions figurant entre le {\mcl begin} et
le {\mcl end} après le {\mcl else} sont toujours exécutées. C'est aussi une erreur que
de mettre un point-virgule `{\idt $;$}' \underline{avant} le {\mcl else} mais elle n'est pas
très grave car elle est signalée par le com\-pi\-la\-teur.

Dans le dessin de la frise en dents de scie ci-dessous, le \robot\ tourne de trois quarts de tour à droite 
s'il est orienté au nord-est et à gauche dans les autres cas (le seul autre cas possible étant d'ailleurs
qu'il soit tourné vers le sud)~:

\vspace{0.4cm}
 {\raggedright 
 \includevspacePostScript[0.5cm]{FriseEnDentsDeScie}{0.95}{0.95}
 \input FriseEnDentsDeScie.tex
 }
 
 Quand il y a plus de deux cas à tester on peut utiliser plusieurs alternatives et ne rien faire
 quand aucune des expressions booléeennes ne vaut {\idt true} comme dans~:
 
\index{Test}
 \begin{center}
\begin{tabular}{@{}|p{12.6cm}|@{}} \hline
{\em Test ``if'' à $n$ alternatives~:}\\ \hline \hline
\begin{minipage}[t]{12.6cm}
{\medskip
\vspace{\neglskip}\mbox{\blc\ \ \ {\mcl if}\hspace{0.3em}\blanc\hspace{0.3em}{\small\rm expression booléenne 1}\hspace{0.3em}\blanc\hspace{0.3em}{\mcl then}}\par 
\vspace{\neglskip}\mbox{\blc\ \ \ \ \mcl begin}\par 
\vspace{\neglskip}\mbox{\blc\ \ \ \ \ \small\rm  \ldots\ Instructions (séparées par des `$;$') exécutées si l'expression }\par 
\vspace{\neglskip}\mbox{\blc\ \ \ \ \ \small\rm  booléenne 1 vaut {\idt true} \ldots}\par 
\vspace{\neglskip}\mbox{\blc\ \ \ \ {\mcl end}}\par 
\vspace{\neglskip}\mbox{\blc\ \ \ {\mcl else}\ {\mcl if}\hspace{0.3em}\blanc\hspace{0.3em}{\small\rm expression booléenne 2}\hspace{0.3em}\blanc\hspace{0.3em}{\mcl then}}\par 
\vspace{\neglskip}\mbox{\blc\ \ \ \ \mcl begin}\par 
\vspace{\neglskip}\mbox{\blc\ \ \ \ \ \small\rm  \ldots\ Instructions (séparées par des `$;$') exécutées si l'expression}\par 
\vspace{\neglskip}\mbox{\blc\ \ \ \ \ \small\rm  booléenne 2 vaut {\idt true} \ldots}\par 
\vspace{\neglskip}\mbox{\blc\ \ \ \ {\mcl end}}\par 
\mbox{\blc\ \ \ \dots}\par
\vspace{\neglskip}\mbox{\blc\ \ \ {\mcl else}\ {\mcl if}\hspace{0.3em}\blanc\hspace{0.3em}{\small\rm expression booléenne $n$}\hspace{0.3em}\blanc\hspace{0.3em}{\mcl then}}\par 
\vspace{\neglskip}\mbox{\blc\ \ \ \ \mcl begin}\par 
\vspace{\neglskip}\mbox{\blc\ \ \ \ \ \small\rm  \ldots\ Instructions (séparées par des `$;$') exécutées si l'expression}\par 
\vspace{\neglskip}\mbox{\blc\ \ \ \ \ \small\rm  booléenne $n$ vaut {\idt true} \ldots}\par 
\vspace{\neglskip}\mbox{\blc\ \ \ \ {\mcl end}\idt$;$}\par 
\medskip}
\end{minipage}\\ \hline 
\end{tabular}
\end{center}

 \noindent On peut aussi prévoir une dernière alternative qui est exécutée quand toutes les
 expressions booléeennes testées valent {\idt false}~:
 
 \index{Test}
\begin{center}
\begin{tabular}{@{}|p{12.6cm}|@{}} \hline
{\em Test ``if'' à $(n + 1)$ alternatives~:}\\ \hline \hline
\begin{minipage}[t]{12.6cm}
{\medskip
\vspace{\neglskip}\mbox{\blc\ \ \ {\mcl if}\hspace{0.3em}\blanc\hspace{0.3em}{\small\rm expression booléenne 1}\hspace{0.3em}\blanc\hspace{0.3em}{\mcl then}}\par 
\vspace{\neglskip}\mbox{\blc\ \ \ \ \mcl begin}\par 
\vspace{\neglskip}\mbox{\blc\ \ \ \ \ \small\rm  \ldots\ Instructions  (séparées par des `$;$') exécutées si l'expression}\par 
\vspace{\neglskip}\mbox{\blc\ \ \ \ \ \small\rm   booléenne 1 vaut {\idt true} \ldots}\par 
\vspace{\neglskip}\mbox{\blc\ \ \ \ {\mcl end}}\par 
\vspace{\neglskip}\mbox{\blc\ \ \ {\mcl else}\ {\mcl if}\hspace{0.3em}\blanc\hspace{0.3em}{\small\rm expression booléenne 2}\hspace{0.3em}\blanc\hspace{0.3em}{\mcl then}}\par 
\vspace{\neglskip}\mbox{\blc\ \ \ \ \mcl begin}\par 
\vspace{\neglskip}\mbox{\blc\ \ \ \ \ \small\rm  \ldots\ Instructions  (séparées par des `$;$') exécutées si l'expression}\par 
\vspace{\neglskip}\mbox{\blc\ \ \ \ \ \small\rm  booléenne 2 vaut {\idt true} \ldots}\par 
\vspace{\neglskip}\mbox{\blc\ \ \ \ {\mcl end}}\par 
\mbox{\blc\ \ \ \dots}\par
\vspace{\neglskip}\mbox{\blc\ \ \ {\mcl else}\ {\mcl if}\hspace{0.3em}\blanc\hspace{0.3em}{\small\rm expression booléenne $n$}\hspace{0.3em}\blanc\hspace{0.3em}{\mcl then}}\par 
\vspace{\neglskip}\mbox{\blc\ \ \ \ \mcl begin}\par 
\vspace{\neglskip}\mbox{\blc\ \ \ \ \ \small\rm  \ldots\ Instructions  (séparées par des `$;$') exécutées si l'expression}\par 
\vspace{\neglskip}\mbox{\blc\ \ \ \ \ \small\rm   booléenne $n$ vaut {\idt true} \ldots}\par 
\vspace{\neglskip}\mbox{\blc\ \ \ \ {\mcl end}}\par 
\vspace{\neglskip}\mbox{\blc\ \ \ {\mcl else}}\par 
\vspace{\neglskip}\mbox{\blc\ \ \ \ \mcl begin}\par 
\vspace{\neglskip}\mbox{\blc\ \ \ \ \ \small\rm  \ldots\ Instructions  (séparées par des `$;$') exécutées si les expressions}\par 
\vspace{\neglskip}\mbox{\blc\ \ \ \ \ \small\rm  booléennes 1, \ldots, $n$ valent {\idt false} \ldots}\par 
\vspace{\neglskip}\mbox{\blc\ \ \ \ {\mcl end}\idt$;$}\par 
\medskip}
\end{minipage}\\ \hline 
\end{tabular}
\end{center}

 Cette forme de test multiple est illustrée ci-dessous dans le programme construisant la frise suivante~:

  {\raggedright 
 \includevspacePostScript[1.0cm]{FriseDeVaguelettes}{0.95}{0.95}
 \input FriseDeVaguelettes.tex
 }

Pour terminer nous allons introduire de nouvelles commandes du \robot\ qui permettent
de faire pivoter le \robot\ sur place vers une orientation donnée. Par exemple la commande
{\idt pn} qui permet de faire \underline{p}ivoter le \robot\ vers le \underline{n}ord peut
être définie à l'aide de la procédure \pascal\ suivante~:

{
\parindent \parindentprog
\medskip

\pagebreak[3] 
\vspace{\neglskip}\mbox{\blc  \blk\idt  \mcl procedure\idt \ pn$;$}\par 
\vspace{\neglskip}\mbox{\blc  \blk\idt  \mcl begin\idt }\par 
\pagebreak[3] 
\vspace{\neglskip}\mbox{\blc  \blk\blk\idt  \mcl if\idt \ tne\ \mcl then\idt }\par 
\vspace{\neglskip}\mbox{\blc  \blk\blk\blk\idt  \mcl begin\idt }\par 
\vspace{\neglskip}\mbox{\blc  \blk\blk\blk\blk\idt  phtg$;$\ }\par 
\vspace{\neglskip}\mbox{\blc  \blk\blk\blk\idt  \mcl end\idt }\par 
\pagebreak[3] 
\vspace{\neglskip}\mbox{\blc  \blk\blk\idt  \mcl else\idt \ \mcl if\idt \ te\ \mcl then\idt }\par 
\vspace{\neglskip}\mbox{\blc  \blk\blk\blk\idt  \mcl begin\idt }\par 
\vspace{\neglskip}\mbox{\blc  \blk\blk\blk\blk\idt  pqtg$;$\ }\par 
\vspace{\neglskip}\mbox{\blc  \blk\blk\blk\idt  \mcl end\idt }\par 
\pagebreak[3] 
\vspace{\neglskip}\mbox{\blc  \blk\blk\idt  \mcl else\idt \ \mcl if\idt \ tse\ \mcl then\idt }\par 
\vspace{\neglskip}\mbox{\blc  \blk\blk\blk\idt  \mcl begin\idt }\par 
\vspace{\neglskip}\mbox{\blc  \blk\blk\blk\blk\idt  p3htg$;$\ }\par 
\vspace{\neglskip}\mbox{\blc  \blk\blk\blk\idt  \mcl end\idt }\par 
\pagebreak[3] 
\vspace{\neglskip}\mbox{\blc  \blk\blk\idt  \mcl else\idt \ \mcl if\idt \ ts\ \mcl then\idt }\par 
\vspace{\neglskip}\mbox{\blc  \blk\blk\blk\idt  \mcl begin\idt }\par 
\vspace{\neglskip}\mbox{\blc  \blk\blk\blk\blk\idt  pdt$;$\ }\par 
\vspace{\neglskip}\mbox{\blc  \blk\blk\blk\idt  \mcl end\idt }\par 
\pagebreak[3] 
\vspace{\neglskip}\mbox{\blc  \blk\blk\idt  \mcl else\idt \ \mcl if\idt \ tsw\ \mcl then\idt }\par 
\vspace{\neglskip}\mbox{\blc  \blk\blk\blk\idt  \mcl begin\idt }\par 
\vspace{\neglskip}\mbox{\blc  \blk\blk\blk\blk\idt  p3htd$;$\ }\par 
\vspace{\neglskip}\mbox{\blc  \blk\blk\blk\idt  \mcl end\idt }\par 
\pagebreak[3] 
\vspace{\neglskip}\mbox{\blc  \blk\blk\idt  \mcl else\idt \ \mcl if\idt \ tw\ \mcl then\idt }\par 
\vspace{\neglskip}\mbox{\blc  \blk\blk\blk\idt  \mcl begin\idt }\par 
\vspace{\neglskip}\mbox{\blc  \blk\blk\blk\blk\idt  pqtd$;$\ }\par 
\vspace{\neglskip}\mbox{\blc  \blk\blk\blk\idt  \mcl end\idt }\par 
\pagebreak[3] 
\vspace{\neglskip}\mbox{\blc  \blk\blk\idt  \mcl else\idt \ \mcl if\idt \ tnw\ \mcl then\idt }\par 
\vspace{\neglskip}\mbox{\blc  \blk\blk\blk\idt  \mcl begin\idt }\par 
\vspace{\neglskip}\mbox{\blc  \blk\blk\blk\blk\idt  phtd$;$\ }\par 
\vspace{\neglskip}\mbox{\blc  \blk\blk\blk\idt  \mcl end\idt $;$\ }\par 
\vspace{\neglskip}\mbox{\blc  \blk\idt  \mcl end\idt $;$}\par 
\medskip
}

\noindent Nous pourrions définir de la même façon les  commandes de pivotement du \robot\
sur place dans l'une des huit directions principales de la rose des vents qui sont données dans la table
\ref{tab:pnetc} ci-dessous.

\begin{table}[htbp]
\begin{center}
\begin{tabular}{@{}|l|p{11.7cm}|@{}} \hline
\multicolumn{2}{|l|}{\em Commandes de pivotement du \robot\ sur place~:} \\ \hline \hline
{\idt pn} & fait \underline{p}ivoter le \robot\ sur place vers le nord (\underline{n}). \\ \hline
{\idt pne} &  fait \underline{p}ivoter le \robot\ sur place  vers le nord-est (\underline{ne}). \\ \hline
{\idt pe} &  fait \underline{p}ivoter le \robot\ sur place vers l'est (\underline{e}). \\ \hline
{\idt pse} &  fait \underline{p}ivoter le \robot\ sur place vers le sud-est (\underline{se}). \\ \hline
{\idt ps} &  fait \underline{p}ivoter le \robot\ sur place vers le sud (\underline{s}). \\ \hline
{\idt psw} &  fait \underline{p}ivoter le \robot\ sur place vers le sud-ouest (\underline{sw}). \\ \hline
{\idt pw} &  fait \underline{p}ivoter le \robot\ sur place vers l' ouest (\underline{w}). \\ \hline
{\idt pnw} &  fait \underline{p}ivoter le \robot\ sur place vers le nord-ouest (\underline{nw}). \\ \hline
\end{tabular}
\end{center}
\caption{Pivotement du \protect\robot\ sur place dans une direction donnée\label{tab:pnetc}}
\end{table}

\noindent \rule{137mm}{0.1mm} %---------------------------------------------
\begin{exercice}{Écrire des programmes \pascal\  comportant des instructions
de test ``if'' pour réaliser les frises de la figure \ref{fig:frisif} ci-dessous.
\noindent $\Box$}\end{exercice}
\begin{figure}[htbp]
 \includevspacePostScript[3.8cm]{FriseIf}{0.95}{0.95}
 \caption{Frises de cercles et (semi-) octogones}\label{fig:frisif}
\end{figure}
\begin{corrige}{
Nous donnons le programme pour dessiner la frise d'octogones de la figure \ref{fig:frisif}
 (page \pageref{fig:frisif}), les autres se trouvant sur disquette~:
\input FriseDoctogones.tex
\noindent $\Box$}\end{corrige}

%%%%%%%%%%%%%%%%%%%%%%%%%%%%%%%%%%%%%%%%%%%%%%%%%%%%%%%%%%%%%%%%%%%%%%%%%%%
\noindent \framebox[\textwidth]{
\begin{minipage}{\myparboxwidth}
\setlength{\parindent}{2em}
\small \sf
\vspace{\abovedisplayskip}
\section*{Addenda}
\index{Test}
L'instruction de test {\idt si} \ldots {\idt alors} \ldots {\idt sinon} \ldots fut introduite vers la fin des années
1950 dans le langage de programmation {\sc algol 60}. Ce langage, conçu par un groupe
d'experts internationaux, eut une très grande influence dans la conception de langages de programmation
ultérieurs notamment {\sc algol w}, \pascal, {\sc ada},  {\sc modula} et bien d'autres.  
Une autre forme d'instruction de test par cas fut introduite dans le langage de programmation 
{\sc algol w} (également conçu par N. Wirth) et reprise en \pascal.  Cette instruction ``case''
généralise l'instruction ``if'' qui peut s'écrire comme suit~:

{
\parindent \parindentprog
\medskip
\vspace{\neglskip}\mbox{\blc  \blk\idt  \mcl case\idt \ Expression\ booléenne\ \mcl of\idt }\par 
\vspace{\neglskip}\mbox{\blc  \blk\blk\idt  true\ \ $:$\ \mcl begin\idt \ \ldots\ \mcl end\idt $;$}\par 
\vspace{\neglskip}\mbox{\blc  \blk\blk\idt  false\ $:$\ \mcl begin\idt \ \ldots\ \mcl end\idt $;$}\par 
\vspace{\neglskip}\mbox{\blc  \blk\idt  \mcl end\idt $;$}\par 
\medskip
 }

De manière plus générale, l'instruction case permet de choisir quelle séquence d'instructions doit
être exécutée en fonction de la valeur d'une expression. Par exemple si {\idt C} est une couleur,
on pourra écrire~:
 
\vspace{\abovedisplayskip}
\end{minipage}
}
%%%%%%%%%%%%%%%%%%%%%%%%%%%%%%%%%%%%%%%%%%%%%%%%%%%%%%%%%%%%%%%%%%%%%%%%%%%% 
%%%%%%%%%%%%%%%%%%%%%%%%%%%%%%%%%%%%%%%%%%%%%%%%%%%%%%%%%%%%%%%%%%%%%%%%%%%
\noindent \framebox[\textwidth]{
\begin{minipage}{\myparboxwidth}
\setlength{\parindent}{2em}
\small \sf
\vspace{\abovedisplayskip}

\index{Test}
{
\parindent \parindentprog
\medskip
\vspace{\neglskip}\mbox{\blc  \blk\idt  \mcl case\idt \ C\ \mcl of\idt }\par 
\vspace{\neglskip}\mbox{\blc  \blk\blk\idt  blanc\ \ $:$\ \mcl begin\idt \ \ldots\ \mcl end\idt $;$}\par 
\vspace{\neglskip}\mbox{\blc  \blk\blk\idt  jaune\ \ $:$\ \mcl begin\idt \ \ldots\ \mcl end\idt $;$}\par 
\vspace{\neglskip}\mbox{\blc  \blk\blk\idt  bleu\ \ \ $:$\ \mcl begin\idt \ \ldots\ \mcl end\idt $;$}\par 
\vspace{\neglskip}\mbox{\blc  \blk\blk\idt  violet\ $:$\ \mcl begin\idt \ \ldots\ \mcl end\idt $;$}\par 
\vspace{\neglskip}\mbox{\blc  \blk\blk\idt  vert\ \ \ $:$\ \mcl begin\idt \ \ldots\ \mcl end\idt $;$}\par 
\vspace{\neglskip}\mbox{\blc  \blk\blk\idt  rouge\ \ $:$\ \mcl begin\idt \ \ldots\ \mcl end\idt $;$}\par 
\vspace{\neglskip}\mbox{\blc  \blk\blk\idt  indigo\ $:$\ \mcl begin\idt \ \ldots\ \mcl end\idt $;$}\par 
\vspace{\neglskip}\mbox{\blc  \blk\blk\idt  noir\ \ \ $:$\ \mcl begin\idt \ \ldots\ \mcl end\idt $;$}\par 
\vspace{\neglskip}\mbox{\blc  \blk\idt  \mcl end\idt $;$}\par 
\medskip
}

\noindent Plusieurs cas semblables peuvent être regroupés comme dans l'exemple ci-dessous~:

{
\parindent \parindentprog
\medskip
\vspace{\neglskip}\mbox{\blc  \blk\idt  \mcl case\idt \ C\ \mcl of\idt }\par 
\vspace{\neglskip}\mbox{\blc  \blk\blk\idt  bleu\ \ \ $:$\ \mcl begin\idt \ \ldots\ \mcl end\idt $;$}\par 
\vspace{\neglskip}\mbox{\blc  \blk\blk\idt  vert\ \ \ $:$\ \mcl begin\idt \ \ldots\ \mcl end\idt $;$}\par 
\vspace{\neglskip}\mbox{\blc  \blk\blk\idt  rouge\ \ $:$\ \mcl begin\idt \ \ldots\ \mcl end\idt $;$}\par 
\vspace{\neglskip}\mbox{\blc  \blk\blk\idt  blanc$,$\ jaune$,$\ violet$,$\ indigo$,$\ noir\ $:$\ \mcl begin\idt \ \ldots\ \mcl end\idt $;$}\par 
\vspace{\neglskip}\mbox{\blc  \blk\idt  \mcl end\idt $;$}\par 
\medskip
}

Tous les cas possibles doivent être prévus. S'il n'y a rien à faire dans certains cas, il faut le prévoir mais
supprimer le {\mcl begin\idt \ \ldots\ \mcl end} correspondant. C'est pourquoi il est prévu dans certaines extensions de 
\pascal\ de pouvoir regrouper tous les cas omis dans une même alternative désignée par le mot
{\mcl otherwise} (`autrement' en anglais) comme dans l'exemple ci-dessous~:

{
\parindent \parindentprog
\medskip
\vspace{\neglskip}\mbox{\blc  \blk\idt  \mcl case\idt \ C\ \mcl of\idt }\par 
\vspace{\neglskip}\mbox{\blc  \blk\blk\idt  bleu\ \ \ \ \ \ $:$\ \mcl begin\idt \ \ldots\ \mcl end\idt $;$}\par 
\vspace{\neglskip}\mbox{\blc  \blk\blk\idt  vert\ \ \ \ \ \ $:$\ \mcl begin\idt \ \ldots\ \mcl end\idt $;$}\par 
\vspace{\neglskip}\mbox{\blc  \blk\blk\idt  rouge\ \ \ \ \ $:$\ \mcl begin\idt \ \ldots\ \mcl end\idt $;$}\par 
\vspace{\neglskip}\mbox{\blc  \blk\blk\idt  \mcl otherwise\idt \ \ \ \mcl begin\idt \ \ldots\ \mcl end\idt $;$}\par 
\vspace{\neglskip}\mbox{\blc  \blk\idt  \mcl end\idt $;$}\par 
\medskip
}

 
\vspace{\abovedisplayskip}
\end{minipage}
}
%%%%%%%%%%%%%%%%%%%%%%%%%%%%%%%%%%%%%%%%%%%%%%%%%%%%%%%%%%%%%%%%%%%%%%%%%%%% 
\newpage                 % Si on termine le chapitre par une page blanche
\thispagestyle{empty}
